% Original vector figure for Version 3.9.0.
\begin{figure}[htbp]
\centering
\begin{tikzpicture}[font=\small,>=Stealth,line width=.65pt]
\filldraw[fill=gray!8,draw=gray] (-2,-1.3)--(.6,-1.3)--(1.8,-.6)--(-.8,-.6)--cycle;
\node[below] at (-.1,-1.3) {$\ker\lambda$};
\fill (0,-.9) circle(1.5pt) node[left] {$0$};
\filldraw[fill=blue!9,draw=blue!70!black] (-2,.2)--(.6,.2)--(1.8,.9)--(-.8,.9)--cycle;
\coordinate (p) at (0,.6);\fill (p) circle(2pt) node[below] {$p$};
\draw[->,blue!80!black,thick] (p)--(1,.6) node[right] {$u$};
\draw[->,blue!80!black,thick] (p)--(.65,.98) node[above right] {$v$};
\draw[->,red!75!black,thick] (p)--(0,2) node[above] {$n$};
\draw[dashed,gray] (0,-.9)--(p);
\node[right,align=left] at (2,.5) {$p+\ker\lambda$\\$\lambda(x)=c$};
\end{tikzpicture}
\caption{An affine plane and its parallel direction plane. The arrows $u,v$ span the direction space; a Euclidean normal $n$ represents the covector. Orthogonality is spatial, not an angle measured on this oblique drawing.}\label{fig:affine-hyperplane}
\end{figure}
